\section{Angelic States of Being}


\label{sec:AngelicStates}

\begin{quotex}
Hear us Lord, holy Lord, almighty Father, eternal God, and deign to send your holy Angel from heaven to guard, cherish, protect, visit and defend all who are gathered together in this place. \flright{\textit{Liturgical prayer of the introductory service preceding the solemn Mass}}

The citizens of heaven are spirits of power, glorious, blessed, having individuality and ranked in order of dignity from their creation, perfect of their kind, with ethereal bodies, immortal, not created incapable of suffering, but made so pure in mind, kindly in affection, devout in piety, perfect in chastity, of one mind in agreement, secure in peace, created by God, dedicated in the praise and worship of God. \flright{\textsc{Saint Bernard of Clairvaux}}

We have gone on in mind, but not with our whole mind; with only part, and that too small a part. Our affections lie weighted down by this bodily mass, and they are stuck to the mire by desire: for now only dry and delicate consideration flies before. \flright{\textsc{Saint Bernard of Clairvaux}}

\end{quotex}
\paragraph{States of Being}
The Person exists in multiple states of being, of which the human state is but one of them. Rene Guenon explains it like this:

\begin{quotex}
almost everything that is said theologically of the angels can also be said metaphysically of the higher states of the being. … the sum of all these states is still nothing at all in relation to the personality, which alone is the true being, because it alone represents its permanent and unconditioned state, and because there is nothing else which can be considered as absolutely real. … in the subtle state we are still only concerned with ideas clothed in forms, since the possibilities which this state comprises do not extend beyond individual existence \flright{\textsc{Rene Guenon}, \emph{Man and his Becoming}}

\end{quotex}
The theologian \textbf{Sergius Bulgakov} describes it in a different way:

\begin{quotex}
everyone has a personal guardian angel who is his heavenly image.

\end{quotex}
So how is this heavenly image related to the human person? He elaborates:

\begin{quotex}
A guardian angel has an affinity of individual character with a human … There exists a likeness between the hypostasis of a guardian angel and that of a human: it is one and the same individuality living in two worlds, in heaven and on earth. … The form of being remains distinct in the spiritual and human world, in heaven and on earth. Here one should not speak about identity but only about correspondence or likeness. Sometimes this idea is expressed in the straightforward convergence of an angel and a human soul which after its liberation from the body assumes a certain luminous shell similar to an angel's. \flright{\textsc{Sergius Bulgakov}, \emph{Jacob's Ladder}}

\end{quotex}
The hypostasis or individuality is what we have been calling the Person, or the I. So the human and spiritual worlds represent different states of the Person. Of course, they are not identical; the guardian angel is not human. Nevertheless, there is necessarily a likeness, otherwise they could not be states of the same person. Specifically, there is a continuity between the human state and angelic states.

Valentin Tomberg, in \emph{Meditations on the Tarot}, goes even further. One of the Hermetic tasks is to ascend Jacob's Ladder by rising through the hierarchy of angels. He explains:

\begin{quotex}
The transcendental Self is not God. It is in his image and after his likeness, according to the law of analogy or kinship, but it is not identical with God. There are still several degrees on the ladder of analogy which separate it from the summit of the ladder —from God. These degrees which are higher than it are its “stars”—or the ideals to which it aims. The Apocalypse specifies the number of them: there are twelve degrees higher than that of the consciousness of the human transcendental Self. It is necessary, therefore, in order to attain to the ONE God, to elevate oneself successively to degrees of consciousness of the nine spiritual hierarchies and the Holy Trinity

\end{quotex}
He even provides a specific example of this ascent in which the Person moves beyond even the Guardian Angel:

\begin{quotex}
The guardian Angel is the friend of the bride at the spiritual marriage of the soul and God. The guardian Angel withdraws before the approach of One who is greater than he. There is what is called in Christian Hermetism the “freeing of the guardian Angel”. The guardian Angel is freed —often in order to be able to acquit new missions—when the soul has acquired the disposition of its part of “likeness” in order to experience the Divine more intimately and more immediately, which corresponds to another hierarchical degree.

Then it is an Archangel who replaces the freed guardian Angel. Human beings whose guardian is an Archangel have not only new experiences of the Divine in their inner life, but also, through this very fact, receive a new and objective vocation. They become representative, of a human group—a nation or a human karmic community—which means to say that from this time onwards their actions will no longer be purely personal but will at the same time have significance and value for those of the human community that they represent.

It also happens sometimes that the Archangel is freed as well. Then it is an entity from the hierarchy of Powers or Elohim which replaces the Archangel. The human being then becomes a representative of the future of humanity. He lives in the present what mankind someday is due to experience in future centuries.

\end{quotex}
\paragraph{Principles and Persons}
For some reason, it is considered more sophisticated to believe in “principles” rather than persons. So there is talk about the “laws of physics”, higher principles, and so on. But do you really regard your wife or husband merely as “principles”, say, for example, of authority or of caring. Are your children “principles”. Do you love them as principles or as the unique persons that they are? Of course, we have our roles to play in the world, and as we act in principled ways, we embody those principles.

Principles or abstract ideas do not have agency; they accomplish nothing without a person with the will to embody those principles. Only a person can make an idea effective in life, the world, in history. Nor does a principle have meaning or value without a person to embody it. This is also true in the spiritual world, as explained by Valentin Tomberg:

\begin{quotex}
The spiritual world is not a world of laws, principles and ideas; it is a world of spiritual beings — human souls, Angels, Archangels, Principalities, Powers, Virtues, Dominions, Thrones, Cherubim, Seraphim and the Holy Trinity: the Holy Spirit, the Son and the Father.

The vertical world, the spiritual world, is that of values and, as the “value of values” is the individual being, it is a world of individual beings or entities.

\end{quotex}
Unlike fallen humans, the angels act with intelligence, justice, and detachment. Here are some examples:

\begin{quotex}
The holy angels punish without anger those whom they receive for punishment by the eternal law of God; they help the suffering without the compassion of pity; and when those whom they love fall in danger, they minister without fear. \flright{\textsc{Saint Augustine}, \emph{City of God}, Book IX}

\end{quotex}
\paragraph{The Body of Angels}
Saint Bernard says that the angels have ethereal bodies, i.e., they are not material. As ethereal, they are not subject to the limitations of space, just to time. That is how they can sometimes be experienced in a manner imitative of sensual experience. After all, sensual experience arises in the soul, it is not “out there”. We experience bodies in our imagination and particularly in dreams. Although dreams unfold in time, there is not space, since the scene can change into another in an instant.

\paragraph{Angelic Hierarchy}
Dionysius first formulated the angelic hierarchy. There are nine orders of angels arranged hierarchically. They each have different roles and functions. These are their roles as described by Saint Bernard.

\subparagraph{Angels}
An Angel is entrusted with the care of a particular person. This angel is the heavenly prototype of that person, and sees God face to face:

\begin{quotex}
in heaven their angels always behold the face of my Father who is in heaven. \flright{\textsc{Matthew, 18:10}}

\end{quotex}
\subparagraph{Archangels}
Archangels know divine mysteries and are sent only on serious and important occasions. The most important occasion was the visit of the archangel Gabriel to Mary.

\subparagraph{Principalities}
The Principalities, through their moderation and wisdom, set up and rule every power on earth which are kept within bounds, transferred, diminished, altered.

\subparagraph{Powers}
The Powers check the powers of darkness and binds the malignity of this air so that it can do no evil, nor any harm, unless it is for our good.

\subparagraph{Virtues}
The Virtues produce the signs and wonders which appear in the elements or are formed from the elements to instruct mortal men. There are many references to these.

\subparagraph{Dominions}
The Dominions minister to spirits of the lower orders. They are masters of the rulership of the principalities, guardians of the powers, the work of the virtues, the revelations of the archangels, the care and prevision of the angels.

\subparagraph{Thrones}
The Thrones are seated and God is seated on them. Sitting is symbolic for supreme tranquillity, placid serenity, the peace which passes understanding.

\subparagraph{Cherubim}
The Cherubim drink from the very fount of Wisdom and pour forth a stream of knowledge to all the citizens of heaven

\subparagraph{Seraphim}
The Seraphim are aflame with the divine fire. They kindle the other citizens so that each is a burning and shining light, and are burning with love, shining with knowledge.

\paragraph{Summary}
These descriptions need to be read on three levels: as descriptions of the angels, as revealing how God acts through the angels, and how they act through us. Bernard reveals how God does different things through different spirits. As for the human element, the angel is suggestive not coercive:

\begin{quotex}
An angel can be present in us. … An angel is “within” when he suggests we do good. He does not enter to cause us to do good. He is exhorting us to do good, not making us good.

\end{quotex}
Saint Bernard summarizes the three aspects like this:

\begin{quotex}
God loves us as Love itself; he knows as Truth itself. He sits as Equity, rules as Majesty, governs as Prince, keeps safe as Salvation, works as Strength, reveals as Light, is with us as Holiness.

\end{quotex}
All these things the angels do.

\begin{quotex}
So do we, but in a far lowlier way. Not because of the good we are but because of the good we share.

\end{quotex}
We can now understand theosis. Just like God, just like the angels:

\emph{We are to love, to know, to be at peace, to rule, to govern, to be strong, to live in the light, to be holy.}

\paragraph{Appendix}
Saint Bernard in \emph{On Consideration} reversed the Principalities and Virtues in the hierarchy, but I restored them to the more traditional ordering of Dionysius. The functions don't change.



\flrightit{Posted on 2020-09-17 by Cologero }
